\documentclass[aspectratio=169,11pt]{beamer}

% --- Theme ---
\usetheme{metropolis}
\usepackage{appendixnumberbeamer}

% --- Packages ---
\usepackage{fontspec}
\usepackage{minted}
\setminted{fontsize=\small,breaklines=true}
\usepackage{booktabs}
\usepackage{tikz}
\usepackage{fontawesome5}

% --- Colors ---
\definecolor{healthgreen}{RGB}{34,139,34}
\definecolor{alertred}{RGB}{220,53,69}
\setbeamercolor{alerted text}{fg=alertred}

% --- Metadata ---
\title{Module 9: Geospatial and temporal health data}
\subtitle{Data analysis with Python for health specialists}
\author{Ya\'{e} Ulrich Gaba}
\institute{AIRINA Labs}
\date{2026}

\begin{document}

% ============================================================
\maketitle

% ============================================================
\begin{frame}{Time series in health}
A health time series has three components:

\begin{itemize}
  \item \textbf{Trend:} long-term direction (incidence rising or falling?)
  \item \textbf{Seasonality:} regular cycles (malaria peaks in rainy season, flu in winter)
  \item \textbf{Noise:} random day-to-day fluctuation
\end{itemize}

\vspace{6pt}
\begin{exampleblock}{Epi curves}
During outbreaks, the epidemic curve is the \alert{most important visualization}. Its shape tells you the transmission pattern: point-source (sharp peak), continuous (plateau), person-to-person (waves).
\end{exampleblock}
\end{frame}

% ============================================================
\begin{frame}[fragile]{Date handling and rolling averages}
\begin{minted}{python}
import pandas as pd

# Parse dates and set as index
df["date"] = pd.to_datetime(df["date"])
df = df.set_index("date")

# Rolling averages (smooth daily noise)
df["cases_7d"] = df["cases"].rolling(window=7).mean()
df["cases_14d"] = df["cases"].rolling(window=14).mean()
\end{minted}

\vspace{4pt}
\textbf{7-day window:} removes day-of-week reporting effects.\\
\textbf{14-day window:} smooths more aggressively for noisy data.
\end{frame}

% ============================================================
\begin{frame}[fragile]{Trend decomposition}
\begin{minted}{python}
from statsmodels.tsa.seasonal import seasonal_decompose

result = seasonal_decompose(df["cases"],
                            model="additive", period=30)

fig, axes = plt.subplots(4, 1, figsize=(12, 10), sharex=True)
result.observed.plot(ax=axes[0], title="Observed")
result.trend.plot(ax=axes[1], title="Trend")
result.seasonal.plot(ax=axes[2], title="Seasonal")
result.resid.plot(ax=axes[3], title="Residual")
\end{minted}

\begin{alertblock}{Period matters}
Daily + monthly cycle $\to$ \texttt{period=30}. Weekly + annual $\to$ \texttt{period=52}. Wrong period = misleading decomposition.
\end{alertblock}
\end{frame}

% ============================================================
\begin{frame}[fragile]{COVID-19: loading JHU data}
\begin{minted}{python}
url = ("https://raw.githubusercontent.com/CSSEGISandData/"
       "COVID-19/master/csse_covid_19_data/"
       "csse_covid_19_time_series/"
       "time_series_covid19_confirmed_global.csv")
confirmed = pd.read_csv(url)

# Filter to 5 African countries
countries = ["South Africa", "Nigeria", "Kenya",
             "Egypt", "Morocco"]
africa = confirmed[confirmed["Country/Region"]\
                   .isin(countries)]

# Reshape wide -> long
africa = africa.groupby("Country/Region")\
               .sum(numeric_only=True)
africa_long = africa.T
africa_long.index = pd.to_datetime(africa_long.index)
\end{minted}
\end{frame}

% ============================================================
\begin{frame}[fragile]{COVID-19: daily cases and 7-day average}
\begin{minted}{python}
# Cumulative -> daily new cases
daily = africa_long.diff().clip(lower=0)

# 7-day rolling average
daily_7d = daily.rolling(7).mean()

fig, ax = plt.subplots(figsize=(14, 6))
for country in countries:
    ax.plot(daily_7d.index, daily_7d[country],
            linewidth=2, label=country)
ax.set_ylabel("Daily new cases (7-day average)")
ax.set_title("COVID-19 in 5 African countries")
ax.legend()
\end{minted}

\texttt{.diff()} converts cumulative to daily. \texttt{.clip(lower=0)} handles data corrections (negative values).
\end{frame}

% ============================================================
\begin{frame}[fragile]{Geospatial data with geopandas}
\begin{minted}{python}
import geopandas as gpd

# Built-in world map
world = gpd.read_file(
    gpd.datasets.get_path("naturalearth_lowres"))
africa_map = world[world["continent"] == "Africa"]

# Choropleth: color regions by a variable
fig, ax = plt.subplots(figsize=(10, 10))
africa_map.plot(column="pop_est", cmap="YlOrRd",
                legend=True, edgecolor="black",
                linewidth=0.5, ax=ax)
ax.set_title("Population in African countries")
ax.set_axis_off()
\end{minted}

\textbf{Install:} \texttt{pip install geopandas mapclassify}
\end{frame}

% ============================================================
\begin{frame}[fragile]{Mapping malaria incidence}
\begin{minted}{python}
# Merge WHO malaria data with Africa map
# (requires matching country names!)
name_map = {
    "United Republic of Tanzania": "Tanzania",
    "Democratic Republic of the Congo": "Dem. Rep. Congo",
}
malaria_latest["country"] = malaria_latest["country"]\
    .replace(name_map)

africa_malaria = africa_map.merge(
    malaria_latest, left_on="name",
    right_on="country", how="left")

africa_malaria.plot(column="incidence_per_1000",
                    cmap="YlOrRd", legend=True,
                    missing_kwds={"color": "lightgrey"})
\end{minted}

\alert{Always check for name mismatches} when merging health data with map data.
\end{frame}

% ============================================================
\begin{frame}{Colormaps for health data}
\begin{table}
\begin{tabular}{lp{8cm}}
\toprule
\textbf{Colormap} & \textbf{Use case} \\
\midrule
\texttt{"YlOrRd"} & Sequential: low-to-high severity (incidence, mortality) \\
\texttt{"RdYlGn\_r"} & Diverging: red = bad, green = good (vaccination coverage) \\
\texttt{"Blues"} & Sequential, neutral (population, counts) \\
\texttt{"RdBu\_r"} & Diverging: correlation matrices \\
\bottomrule
\end{tabular}
\end{table}

\vspace{6pt}
\alert{Avoid} rainbow colormaps (\texttt{"jet"}) --- they mislead the eye and are not colorblind-safe.
\end{frame}

% ============================================================
\begin{frame}[fragile]{Small multiples: maps over time}
\begin{minted}{python}
dates_to_plot = ["2020-06-01", "2021-01-01", "2021-06-01"]

fig, axes = plt.subplots(1, 3, figsize=(18, 7))
for ax, date_str in zip(axes, dates_to_plot):
    # Merge case data for this date with map
    temp = africa_map.copy()
    # ... merge logic ...
    temp.plot(column="cases", cmap="Reds", ax=ax,
              edgecolor="black", linewidth=0.5)
    ax.set_title(date_str)
    ax.set_axis_off()

plt.suptitle("COVID-19 cumulative cases in Africa")
\end{minted}

Small multiples are more \textbf{scientifically rigorous} than animations for publications.
\end{frame}

% ============================================================
\begin{frame}{What you can do after this module}
\begin{enumerate}
  \item Parse dates, build datetime indices, and compute rolling averages
  \item Decompose time series into trend, seasonality, and noise
  \item Load and reshape COVID-19 data for epidemic curve analysis
  \item Create choropleth maps of disease burden with geopandas
  \item Merge health data with geospatial boundaries
  \item Build multi-panel dashboards combining time and space
\end{enumerate}

\vspace{12pt}
\centering
\Large \alert{Next: Module 10 --- Capstone project}
\end{frame}

% ============================================================
\begin{frame}{Exercises (Hour 3)}
\begin{enumerate}
  \item \textbf{Deaths curve:} Plot COVID-19 deaths (7-day avg) for 5 African countries
  \item \textbf{Choropleth:} Map COVID-19 cases per million across Africa
  \item \textbf{Small multiples:} Three maps showing cumulative cases at 3 time points
  \item \textbf{Mini-project:} Build a 4-panel COVID-19 dashboard (cases, deaths, vaccination, choropleth)
\end{enumerate}
\end{frame}

\end{document}
